\documentclass{article}
\usepackage{amsmath,amssymb,hyperref}
\begin{document}
\section*{Sharp rational orbit bounds for quadratic rational maps}
\subsection{Definitions and mathematical statement}
Let $f:\mathbb P^1\to\mathbb P^1$ be a degree-two rational map defined over
$\mathbb Q$. It has a presentation $f([X:Y])=[F(X,Y):G(X,Y)]$ by homogeneous
quadratic forms over $\mathbb Q$ with no common projective zero.
The rational projective line includes infinity.

For $P\in\mathbb P^1(\mathbb Q)$ put
$\mathcal O_f(P)=\{f^n(P):n\geq0\}$. A point is preperiodic if this orbit
is finite, and periodic if $f^n(P)=P$ for some positive $n$; the least such
$n$ is its exact period. Let $\operatorname{PrePer}(f,\mathbb Q)$ denote
the set of rational preperiodic points.

Do all three bounds hold for every $f$?
\begin{enumerate}
\item Every rational periodic point has exact period at most $7$.
\item Every rational preperiodic point has at most $9$ distinct points in its
forward orbit, including itself.
\item $\#\operatorname{PrePer}(f,\mathbb Q)\leq14$.
\end{enumerate}

\subsection{Short English statement}
For quadratic rational maps over the rationals, are the universal bounds
$7$ on periods, $9$ on finite orbit sizes, and $14$ on all preperiodic points
valid?

\subsection{Sources}
R. Benedetto, P. Ingram, R. Jones, M. Manes, J. H. Silverman
and T. J. Tucker, \emph{Current Trends and Open Problems in Arithmetic
Dynamics}, arXiv:1806.04980v2 (2018), Conjecture 4.7,
\url{https://arxiv.org/html/1806.04980v2}.
The survey attributes the conjecture to Benedetto, Chen, Hyde, Kovacheva
and White (2014), and Blanc, Canci and Elkies (2015).
This is an independently written formulation of the three assertions.
\end{document}
